\itemhead{习题3.2 · 第12题（教材p.86）}{XJ-B1-C03-HW-P086-Y1121}
\label{XJ-B1-C03-HW-P086-Y1121}
{\footnotesize 湘教版教材p.86 / PDF第93页 / 习题3.2 / 第12题}
\field{原题}{已知 $f(x),g(x)$ 分别是定义在 $\mathbb R$ 上的偶函数和奇函数，且 $f(x)-g(x)=x^3+x^2+1$，求 $f(1)+g(1)$。}
\field{参考答案}{将 $x$ 换成 $-x$ 后相加相减得 $f(x)=x^2+1,g(x)=-x^3$，故 $f(1)+g(1)=1$。}
{\small \field{关联知识}{F08 奇偶性、对称与函数性质应用（教材帮p.113、114、115 / PDF122、123、124）；F03 函数求值、解析式与函数建模（教材帮p.80、87、88、89、90、131 / PDF89、96、97、98、99、140）}
\field{方法与易错点}{检查定义域，再代入 $-x$；求另一侧解析式时同时写出 $-x$ 所在区间。 易错：只看公式不看定义域；用一对点证明奇偶；负半轴代换遗漏最外层负号。}
\field{答案依据}{答案册相关题组：PDF第23、24页（册内页22、23）；按当前题库原题逐项重算，不按旧版题号机械对套。}}
